%% ──────────────────────────────────────────────
%%  پیوست ت – جدول تطبیقی جامع
%%  فایل: appendices/app-d-comparative-table.tex
%% ──────────────────────────────────────────────
\chapter{جدول تطبیقی جامع: ۳۰ مورد مداخله‌ی خارجی}
\label{app:comparative-table}

\begin{tcolorbox}[mybox, title={درباره‌ی این پیوست}]
این پیوست شامل جدول تطبیقی جامعی از بیش از ۳۰ مورد مداخله‌ی خارجی بررسی‌شده در این کتاب است. ستون‌ها شامل: مورد، سال، مداخله‌گر، نوع مداخله، آیا دعوت شده؟، آیا مبنای حقوقی داشت؟، نتیجه‌ی فوری، نتیجه‌ی بلندمدت، و جایگاه بر طیف A تا F. جدول به‌صورت \lat{landscape} و با رنگ‌بندی بر اساس طیف ارائه شده است.
\end{tcolorbox}

\begin{landscape}
\begin{longtable}{
    >{\raggedleft\arraybackslash}p{2.5cm}
    >{\centering\arraybackslash}p{1.2cm}
    >{\raggedleft\arraybackslash}p{2.2cm}
    >{\raggedleft\arraybackslash}p{1.8cm}
    >{\centering\arraybackslash}p{1cm}
    >{\centering\arraybackslash}p{1cm}
    >{\raggedleft\arraybackslash}p{2.5cm}
    >{\raggedleft\arraybackslash}p{3cm}
    >{\centering\arraybackslash}p{1cm}
}

\caption{جدول تطبیقی جامع: مداخلات خارجی از ایران باستان تا ۲۰۲۴} \label{tab:app-d-master} \\

\toprule
\rowcolor{PrimaryDeep}
\textcolor{white}{\textbf{مورد}} &
\textcolor{white}{\textbf{سال}} &
\textcolor{white}{\textbf{مداخله‌گر}} &
\textcolor{white}{\textbf{نوع}} &
\textcolor{white}{\textbf{دعوت؟}} &
\textcolor{white}{\textbf{حقوقی؟}} &
\textcolor{white}{\textbf{نتیجه‌ی فوری}} &
\textcolor{white}{\textbf{نتیجه‌ی بلندمدت}} &
\textcolor{white}{\textbf{طیف}} \\
\midrule
\endfirsthead

\multicolumn{9}{c}{\tablename\ \thetable{} — ادامه} \\
\toprule
\rowcolor{PrimaryDeep}
\textcolor{white}{\textbf{مورد}} &
\textcolor{white}{\textbf{سال}} &
\textcolor{white}{\textbf{مداخله‌گر}} &
\textcolor{white}{\textbf{نوع}} &
\textcolor{white}{\textbf{دعوت؟}} &
\textcolor{white}{\textbf{حقوقی؟}} &
\textcolor{white}{\textbf{نتیجه‌ی فوری}} &
\textcolor{white}{\textbf{نتیجه‌ی بلندمدت}} &
\textcolor{white}{\textbf{طیف}} \\
\midrule
\endhead

\midrule
\multicolumn{9}{r}{\footnotesize ادامه در صفحه‌ی بعد...} \\
\endfoot

\bottomrule
\endlastfoot

%% ── ایران باستان و میانه ──
حمله‌ی اسکندر به ایران & ۳۳۰ ق.م & مقدونیه & فتح نظامی & خیر & — & سقوط هخامنشیان & هلنیزاسیون + احیای ایرانی & C \\
\rowcolor{NeutralLight}
فتح عرب ایران & ۶۵۱ & خلافت اسلامی & فتح & خیر & — & سقوط ساسانیان & اسلامی‌شدن + حفظ زبان/فرهنگ & C \\
حمله‌ی مغول & ۱۲۱۹ & مغول & فتح ویرانگر & خیر & — & ویرانی عظیم & احیای تدریجی (ایلخانان) & E→C \\
\rowcolor{NeutralLight}
حمله‌ی تیمور & ۱۳۸۳ & تیمور & فتح & خیر & — & ویرانی مجدد & احیا تحت ترکمانان & D \\

%% ── صفوی تا قاجار ──
جنگ ایران و روسیه (اول) & ۱۸۰۴ & روسیه & جنگ & خیر & — & گلستان & از دست دادن قفقاز شمالی & D \\
\rowcolor{NeutralLight}
جنگ ایران و روسیه (دوم) & ۱۸۲۶ & روسیه & جنگ & خیر & — & ترکمانچای & از دست دادن قفقاز جنوبی & D \\
امتیاز رویتر & ۱۸۷۲ & بریتانیا & اقتصادی & بلی & — & واگذاری منابع & لغو تحت فشار روسیه & D \\
\rowcolor{NeutralLight}
قرارداد ۱۹۰۷ & ۱۹۰۷ & بریتانیا/روسیه & سیاسی & خیر & — & تقسیم ایران به مناطق نفوذ & تضعیف حاکمیت & D \\

%% ── ایران معاصر ──
اشغال ایران & ۱۹۴۱ & متفقین & اشغال & خیر & — & خلع رضاشاه & وابستگی سیاسی & D \\
\rowcolor{NeutralLight}
بحران آذربایجان & ۱۹۴۶ & شوروی & جدایی‌طلبی & خیر & خیر & جمهوری خودمختار & سرکوب و بازگشت & C \\
کودتای ۲۸ مرداد & ۱۹۵۳ & آمریکا/بریتانیا & کودتا & خیر & خیر & سقوط مصدق & انقلاب ۱۹۷۹ & D→انقلاب \\

%% ── اروپا ──
\rowcolor{NeutralLight}
انقلاب شکوهمند & ۱۶۸۸ & هلند (ویلیام) & دعوت‌شده & بلی & تا حدی & تغییر رژیم & دموکراسی پارلمانی & A \\
ناپلئون در اروپا & ۱۷۹۹– & فرانسه & فتح & خیر & خیر & سلطه‌ی موقت & کنگره‌ی وین & C→B \\
\rowcolor{NeutralLight}
کنگره‌ی وین & ۱۸۱۵ & ائتلاف اروپایی & دیپلماتیک & — & — & بازگشت نظم & صلح نسبی ۱۰۰ ساله & B \\

%% ── استعمار ──
تقسیم آفریقا & ۱۸۸۴ & قدرت‌های اروپایی & استعمار & خیر & خیر & مرزهای مصنوعی & بحران‌های مداوم & F \\
\rowcolor{NeutralLight}
استعمار هند & ۱۷۵۷– & بریتانیا & استعمار & خیر & خیر & شرکت هند شرقی & استقلال ۱۹۴۷ + تجزیه & F→B \\
استعمار الجزایر & ۱۸۳۰ & فرانسه & استعمار & خیر & خیر & استعمار اشغالی & استقلال ۱۹۶۲ پس از جنگ خونین & F→B \\
\rowcolor{NeutralLight}
استعمار کنگو & ۱۹۰۸ & بلژیک & استعمار & خیر & خیر & غارت و بردگی & لومومبا، موبوتو، هرج‌ومرج & F→E \\

%% ── مداخلات معاصر ──
جنگ خلیج فارس & ۱۹۹۱ & ائتلاف ۳۴ کشور & نظامی & بلی & بلی & آزادی کویت & تحریم‌ها، مناطق ممنوع & B \\
\rowcolor{NeutralLight}
سومالی & ۱۹۹۲ & آمریکا/سازمان ملل & انسان‌دوستانه & خیر & بلی & شکست & هرج‌ومرج مداوم & E \\
رواندا & ۱۹۹۴ & — (عدم مداخله) & عدم مداخله & — & — & نسل‌کشی & ثبات اقتدارگرا & خاص \\
\rowcolor{NeutralLight}
کوزوو & ۱۹۹۹ & ناتو & بمباران & خیر & خیر & استقلال & وابستگی به KFOR & B/C \\
عراق & ۲۰۰۳ & آمریکا/بریتانیا & اشغال & خیر & خیر & سقوط صدام & جنگ فرقه‌ای، داعش & C/D \\
\rowcolor{NeutralLight}
افغانستان & ۲۰۰۱ & آمریکا/ناتو & اشغال & تا حدی & تا حدی & سقوط طالبان & بازگشت طالبان ۲۰۲۱ & E \\
لیبی & ۲۰۱۱ & ناتو & بمباران (R2P) & تا حدی & بلی & سقوط قذافی & هرج‌ومرج & E \\
\rowcolor{NeutralLight}
سوریه & ۲۰۱۱– & چندجانبه & نیابتی/مستقیم & متفاوت & خیر & جنگ داخلی & فاجعه‌ی انسانی & D/E \\
یمن & ۲۰۱۵– & عربستان/ائتلاف & بمباران & تا حدی & خیر & جنگ فرسایشی & بحران انسانی & E \\
\rowcolor{NeutralLight}
اوکراین (کریمه) & ۲۰۱۴ & روسیه & الحاق & خیر & خیر & الحاق کریمه & تنش مداوم & F \\
اوکراین (تهاجم) & ۲۰۲۲– & روسیه/غرب & تهاجم/حمایت & خیر/بلی & خیر & مقاومت & ادامه‌دار & ؟ \\
\rowcolor{NeutralLight}
کنگو (جنگ‌ها) & ۱۹۹۶– & رواندا/اوگاندا & منطقه‌ای & خیر & خیر & سقوط موبوتو & ۵ میلیون تلفات & E \\

%% ── آمریکای لاتین ──
کودتای گواتمالا & ۱۹۵۴ & آمریکا (سیا) & کودتا & خیر & خیر & سقوط آربنز & ۳۶ سال جنگ داخلی & D \\
\rowcolor{NeutralLight}
کودتای شیلی & ۱۹۷۳ & آمریکا (سیا) & کودتا & خیر & خیر & سقوط آلنده & دیکتاتوری پینوشه & D \\

\end{longtable}
\end{landscape}

\needspace{6\baselineskip}
\section{تحلیل آماری مختصر}
\label{sec:app-d-analysis}

از ۳۲ مورد فهرست‌شده در جدول فوق:

\begin{tcolorbox}[wavebox, title={خلاصه‌ی آماری ۳۲ مورد مداخله}]
\begin{itemize}
\item \textbf{طیف A (رهایی کامل):} ۱ مورد (انقلاب شکوهمند ۱۶۸۸) = \textbf{۳٪}
\item \textbf{طیف B (رهایی با وابستگی):} ۴ مورد (کویت ۱۹۹۱، وین ۱۸۱۵، و دو مورد پسااستعماری) = \textbf{۱۳٪}
\item \textbf{طیف C (تغییر رژیم + بی‌ثباتی):} ۶ مورد = \textbf{۱۹٪}
\item \textbf{طیف D (وابستگی ساختاری):} ۹ مورد = \textbf{۲۸٪}
\item \textbf{طیف E (هرج‌ومرج):} ۸ مورد = \textbf{۲۵٪}
\item \textbf{طیف F (استعمار/الحاق):} ۴ مورد = \textbf{۱۳٪}
\end{itemize}

\textbf{نتیجه:} تنها \textbf{۱۶٪} از مداخلات (۵ از ۳۲ مورد) نتیجه‌ی «نسبتاً مثبت» (A یا B) داشته‌اند.\\
\textbf{۸۴٪} مداخلات نتایجی از بی‌ثباتی تا فاجعه به‌بار آورده‌اند.
\end{tcolorbox}

\begin{tcolorbox}[quotebox, title={یک عدد که همه‌چیز را خلاصه می‌کند}]
\begin{center}
\Huge\textbf{۸۴٪}\\[8pt]
\large\rl{از مداخلات خارجی بررسی‌شده در این کتاب}\\[4pt]
\large\rl{نتایجی از بی‌ثباتی تا فاجعه به‌بار آورده‌اند.}
\end{center}
\end{tcolorbox}

\cleardoublepage